\chapter{共同结构：Noether、线性响应与规范}

\begin{mapbox}
本章不重复完整课程，而集中打通四条面试常用链路：
\[
\begin{aligned}
\text{正则结构}&\to\text{量子对易关系},\\
\text{小振动}&\to\text{声子与磁振子},
\text{规范不变性}&\to\text{超导相位与 Berry 连接},\\
\text{介质响应}&\to\text{光学、电导与集体模}.
\end{aligned}
\]
\end{mapbox}

\section{作用量、边界项与 Noether 定理}
作用量
$S[q]=\int_{t_1}^{t_2}L(q,\dot q,t)\dd t$ 的变分为
\[
\delta S=\int_{t_1}^{t_2}
\left(\pdv Lq-\dv{}{t}\pdv L{\dot q}\right)\delta q\,\dd t
+\left.p\,\delta q\right|_{t_1}^{t_2}.
\]
固定端点消去边界项，得到 Euler--Lagrange 方程。若变换
$q_i\to q_i+\epsilon\Delta q_i$ 使 Lagrangian 只改变全导数
$\delta L=\epsilon\dd F/\dd t$，则
\[
\boxed{Q=\sum_i p_i\Delta q_i-F}
\]
守恒。这一形式说明“作用量不变”可允许 Lagrangian 改变边界项。

\begin{examplebox}
\textbf{时间平移与能量。}令 $\delta t=\epsilon$，沿轨道
$\delta q_i=\dot q_i\epsilon$。若 $L$ 不显含时间，则守恒量
\[
E=\sum_i p_i\dot q_i-L=H.
\]
若外参量显含时间，系统与驱动器交换能量，机械系统自身的 $H$ 不再守恒。
\end{examplebox}

\section{正则结构与量子化}
Hamilton 方程可写成
\[
\dot z^a=J^{ab}\partial_bH,\qquad
J=\begin{pmatrix}0&I\\-I&0\end{pmatrix},
\quad z=(q,p).
\]
正则变换保持辛矩阵 $J$，也保持 Poisson 括号。量子化的对应关系
\[
\{A,B\}_{\rm PB}\longrightarrow\frac1{i\hbar}[\hat A,\hat B]
\]
在基本正则变量和低阶多项式上给出直觉，但并非任意经典函数都能无歧义量子化。

相空间流散度为
\[
\nabla_z\cdot\dot{\vb z}
=\sum_i\left(\pdv{\dot q_i}{q_i}+
\pdv{\dot p_i}{p_i}\right)=0,
\]
所以 Hamilton 流保持相体积。微正则系综沿轨道不变，Liouville 定理由此成为经典统计力学的几何基础。

\section{小振动、Goldstone 模与软模}
平衡点附近 Hessian $K$ 的零本征值表示没有二次恢复力。若零模来自连续对称性自发破缺，其长波推广就是 Goldstone 模。例如晶体破坏连续平移，对应低能声学声子；各向同性交换铁磁破坏自旋旋转，对应磁振子。

若调参使某个有限波矢模式频率
$\omega^2(\vb Q)$ 下降到零，称软模。软模凝聚可产生结构相变或密度波。动力学稳定要求所有物理模式
$\omega^2\ge0$；虚频意味着所选结构不是局域极小值。

\section{Maxwell 边界条件}
跨越面电荷 $\sigma_f$ 与面电流 $\vb K_f$ 的界面，取薄 pillbox 和小矩形回路：
\[
\hat{\vb n}\cdot(\vb D_2-\vb D_1)=\sigma_f,\qquad
\hat{\vb n}\cdot(\vb B_2-\vb B_1)=0,
\]
\[
\hat{\vb n}\times(\vb E_2-\vb E_1)=0,\qquad
\hat{\vb n}\times(\vb H_2-\vb H_1)=\vb K_f.
\]
它们是反射、折射、表面等离激元和波导模式的起点。若材料参数突变但无自由面源，切向 $E,H$ 连续，法向 $D,B$ 连续。

\section{介电函数、光学电导与 Drude 响应}
频域中取 $e^{-i\omega t}$ 约定，
\[
\vb J(\omega)=\sigma(\omega)\vb E(\omega),\qquad
\varepsilon(\omega)=1+\frac{i\sigma(\omega)}
{\varepsilon_0\omega}.
\]
Drude 方程
$m(\dot{\vb v}+\vb v/\tau)=-e\vb E_0e^{-i\omega t}$ 给
\[
\boxed{\sigma(\omega)=\frac{ne^2\tau/m}{1-i\omega\tau}},
\]
\[
\mathrm{Re}\,\sigma=\frac{ne^2\tau/m}{1+\omega^2\tau^2},
\qquad
\mathrm{Im}\,\sigma=\frac{ne^2\tau/m\,\omega\tau}
{1+\omega^2\tau^2}.
\]
无阻尼极限下
$\varepsilon(\omega)=1-\omega_p^2/\omega^2$。当 $\omega<\omega_p$，介电函数实部为负，体内波矢主要为虚数，金属强反射；当 $\varepsilon(\omega)=0$，出现纵向等离激元。

\begin{examplebox}
\textbf{由光学谱读散射率。}Drude 峰半宽尺度为 $1/\tau$，谱权重尺度为 $ne^2/m^*$。降温后若峰变窄而总低频权重近似保持，说明散射率下降；若权重转移到有限频率，可能涉及能隙、局域化或关联重整化。
\end{examplebox}

\section{Kramers--Kronig 与因果性}
因果响应 $\chi(t<0)=0$ 使 $\chi(\omega)$ 在复上半平面解析，实虚部满足
\[
\mathrm{Re}\,\chi(\omega)=\frac1\pi\mathcal P
\int_{-\infty}^{\infty}
\frac{\mathrm{Im}\,\chi(\omega')}{\omega'-\omega}\dd\omega',
\]
\[
\mathrm{Im}\,\chi(\omega)=-\frac1\pi\mathcal P
\int_{-\infty}^{\infty}
\frac{\mathrm{Re}\,\chi(\omega')}{\omega'-\omega}\dd\omega'.
\]
吸收与色散不是独立的。实验只测有限频段时必须外推，因此 Kramers--Kronig 反演含模型依赖。

\section{规范、Aharonov--Bohm 与超导相位}
带电粒子的最小耦合
\[
H=\frac1{2m}(\vb p-q\vb A)^2+q\phi.
\]
规范变换势的同时，波函数必须变为
\[
\psi'=\exp\left(\frac{iq\chi}{\hbar}\right)\psi,
\]
协变导数才保持物理不变。即使粒子路径处 $\vb B=0$，闭合回路相位
\[
\Delta\varphi=\frac q\hbar\oint\vb A\cdot\dd\vb l
=\frac q\hbar\Phi
\]
可由磁通决定，这是 Aharonov--Bohm 效应。

超导序参量
$\Psi=|\Psi|e^{i\theta}$ 的规范不变量是
\[
\hbar\nabla\theta-q\vb A.
\]
闭合回路单值性给
$\oint\nabla\theta\cdot\dd\vb l=2\pi n$；深处超流速度近零时得到磁通量子
\[
\Phi_0=\frac h{2e}.
\]
分母 $2e$ 直接反映 Cooper 对电荷。

\section{三个完整例题}
\begin{examplebox}
\textbf{例 1：两耦合振子的声学/光学类比。}相同质量、胞内外弹簧不同的二自由度系统，改用
$Q_\pm=(q_1\pm q_2)/\sqrt2$。同相坐标主要对应整体移动，反相坐标拉伸胞内弹簧。这个坐标变换就是正规模分解，也是双原子晶格两类分支的最简原型。
\end{examplebox}
\begin{examplebox}
\textbf{例 2：金属反射边。}忽略阻尼时
$n^2(\omega)=\varepsilon=1-\omega_p^2/\omega^2$。低于 $\omega_p$ 时折射率主要为虚数，场指数衰减；高于 $\omega_p$ 可传播。真实材料还叠加带间跃迁和晶格背景介电常数。
\end{examplebox}
\begin{examplebox}
\textbf{例 3：Josephson 频率。}两超导体相位差 $\delta$ 满足
$\hbar\dot\delta=2eV$，故恒定电压下
\[
I=I_c\sin\delta(t),\qquad f=\frac{2e}{h}V.
\]
它把电压精确转换为频率，展示宏观相位的可观测性。
\end{examplebox}

\clearpage
\begin{selftestbox}
\small
\begin{enumerate}
\item Lagrangian 改变全导数为何不改变运动方程？
\item Liouville 定理与微正则系综有什么关系？
\item 零模、软模与不稳定虚频如何区分？
\item 从 Maxwell 方程推导四个界面条件。
\item $\sigma(\omega)$ 与 $\varepsilon(\omega)$ 如何互换？
\item Drude 峰宽和谱权重各测什么？
\item Kramers--Kronig 关系为何来自因果性？
\item 波函数为什么必须随规范势一起变换？
\item 磁通量子为什么是 $h/2e$？
\item 规范对称与自发对称性破缺的表述为何需要谨慎？
\end{enumerate}
\end{selftestbox}
